% GENERATED FILE. Edit canonical lessons, then run npm run generate:books.
% canonical-source-hash: fnv1a64:35018dcd5f92b612
% canonical-lessons: HI-C518-tar, HI-C518-balb, HI-C518-svic, HI-C518-khuti, HI-C518-tagna

\chapter{Wires and Switches}
\label{ch:hi-wires-and-switches}
\hlchaptermodality{518}

\begin{chapteropening}
\textbf{By the end of this chapter:} \emph{I can talk about wires, bulbs and switches.}

Complete the last lesson of chapter 518: I can talk about wires, bulbs and switches.
\end{chapteropening}

\section[\texorpdfstring{\hi{तार} (tār)}{तार (tār)}]{\hi{तार} (tār) — a wire; a string}

\label{lesson:HI-C518-tar}



\begin{quote}
Before the new one: say the Hindi for to smear (a floor or wall with mud), then the Hindi for to lay (bricks).
\end{quote}

\subsection*{You'll want to know: \hi{तार}}
\textbf{\hi{तार}} — \emph{tār} — \textquotedblleft{}a wire; a string\textquotedblright{}.

\begin{quote}
\hi{बिजली} \hi{का} \hi{तार} \hi{बहुत} \hi{पुराना} \hi{है।}
\end{quote}

\begin{quote}
\emph{Bijlī kā tār bahut purānā hai.}
\end{quote}

\begin{quote}
The electric wire is very old.
\end{quote}

Masculine.

\begin{cousinweb}
\textbf{\hi{तार}} (\emph{tār}) is usually traced to Persian \emph{tār}, \textquotedblleft{}a thread, a string\textquotedblright{}.
\end{cousinweb}

\begin{grammarlens}[title={What you've built}]
One more word for describing what you lived and felt.
\end{grammarlens}

\subsection*{Your turn}
\begin{itemize}
\raggedright
  \item \emph{Say it:} \emph{tār}
  \item \emph{Say it:} \emph{tār}, once more
  \item \emph{Say it:} \emph{Bijlī kā tār bahut purānā hai.} — The electric wire is very old.
  \item \emph{Recall:} say the Hindi for to smear (a floor or wall with mud), then the Hindi for to lay (bricks), then say \emph{tār} again
\end{itemize}

\begin{tcolorbox}[breakable,colback=teal!4,colframe=teal!35!black,title={Before you move on}]
What does \hi{तार} mean? (\textquotedblleft{}A wire; a string\textquotedblright{}.) Say it once more.
\end{tcolorbox}
\section[\texorpdfstring{\hi{बल्ब} (balb)}{बल्ब (balb)}]{\hi{बल्ब} (balb) — a light bulb}

\label{lesson:HI-C518-balb}



\begin{quote}
Before the new one: say the Hindi for to lay (bricks), then the Hindi for a wire.
\end{quote}

\subsection*{You'll want to know: \hi{बल्ब}}
\textbf{\hi{बल्ब}} — \emph{balb} — \textquotedblleft{}a light bulb\textquotedblright{}.

\begin{quote}
\hi{रसोई} \hi{का} \hi{बल्ब} \hi{नहीं} \hi{जल} \hi{रहा} \hi{है।}
\end{quote}

\begin{quote}
\emph{Rasoī kā balb nahīṃ jal rahā hai.}
\end{quote}

\begin{quote}
The kitchen bulb is not lighting up.
\end{quote}

Masculine.

\begin{cousinweb}
\textbf{\hi{बल्ब}} (\emph{balb}) is a loan from English \emph{bulb}.
\end{cousinweb}

\begin{grammarlens}[title={What you've built}]
One more word for describing what you lived and felt.
\end{grammarlens}

\subsection*{Your turn}
\begin{itemize}
\raggedright
  \item \emph{Say it:} \emph{balb}
  \item \emph{Say it:} \emph{balb}, once more
  \item \emph{Say it:} \emph{Rasoī kā balb nahīṃ jal rahā hai.} — The kitchen bulb is not lighting up.
  \item \emph{Recall:} say the Hindi for to lay (bricks), then the Hindi for a wire, then say \emph{balb} again
\end{itemize}

\begin{tcolorbox}[breakable,colback=teal!4,colframe=teal!35!black,title={Before you move on}]
What does \hi{बल्ब} mean? (\textquotedblleft{}A light bulb\textquotedblright{}.) Say it once more.
\end{tcolorbox}
\section[\texorpdfstring{\hi{स्विच} (svic)}{स्विच (svic)}]{\hi{स्विच} (svic) — a switch}

\label{lesson:HI-C518-svic}



\begin{quote}
Before the new one: say the Hindi for a wire, then the Hindi for a light bulb.
\end{quote}

\subsection*{You'll want to know: \hi{स्विच}}
\textbf{\hi{स्विच}} — \emph{svic} — \textquotedblleft{}a switch\textquotedblright{}.

\begin{quote}
\hi{स्विच} \hi{दरवाज़े} \hi{के} \hi{पास} \hi{है।}
\end{quote}

\begin{quote}
\emph{Svic darvāze ke pās hai.}
\end{quote}

\begin{quote}
The switch is by the door.
\end{quote}

Masculine.

\begin{cousinweb}
\textbf{\hi{स्विच}} (\emph{svic}) is a loan from English \emph{switch}.
\end{cousinweb}

\begin{grammarlens}[title={What you've built}]
One more word for describing what you lived and felt.
\end{grammarlens}

\subsection*{Your turn}
\begin{itemize}
\raggedright
  \item \emph{Say it:} \emph{svic}
  \item \emph{Say it:} \emph{svic}, once more
  \item \emph{Say it:} \emph{Svic darvāze ke pās hai.} — The switch is by the door.
  \item \emph{Recall:} say the Hindi for a wire, then the Hindi for a light bulb, then say \emph{svic} again
\end{itemize}

\begin{tcolorbox}[breakable,colback=teal!4,colframe=teal!35!black,title={Before you move on}]
What does \hi{स्विच} mean? (\textquotedblleft{}A switch\textquotedblright{}.) Say it once more.
\end{tcolorbox}
\section[\texorpdfstring{\hi{खूँटी} (khū̃ṭī)}{खूँटी (khū̃ṭī)}]{\hi{खूँटी} (khū̃ṭī) — a peg; a hook (on a wall)}

\label{lesson:HI-C518-khuti}



\begin{quote}
Before the new one: say the Hindi for a light bulb, then the Hindi for a switch.
\end{quote}

\subsection*{You'll want to know: \hi{खूँटी}}
\textbf{\hi{खूँटी}} — \emph{khū̃ṭī} — \textquotedblleft{}a peg; a hook (on a wall)\textquotedblright{}.

\begin{quote}
\hi{दरवाज़े} \hi{के} \hi{पीछे} \hi{एक} \hi{खूँटी} \hi{है।}
\end{quote}

\begin{quote}
\emph{Darvāze ke pīche ek khū̃ṭī hai.}
\end{quote}

\begin{quote}
There is a peg behind the door.
\end{quote}

Feminine.

\begin{cousinweb}
\textbf{\hi{खूँटी}} (\emph{khū̃ṭī}) is a tadbhava, inherited from \emph{*khuṭṭa}, \textquotedblleft{}a peg, a post\textquotedblright{}, a form CDIAL reconstructs from its descendants (CDIAL 3893).
\end{cousinweb}

\begin{grammarlens}[title={What you've built}]
One more word for describing what you lived and felt.
\end{grammarlens}

\subsection*{Your turn}
\begin{itemize}
\raggedright
  \item \emph{Say it:} \emph{khū̃ṭī}
  \item \emph{Say it:} \emph{khū̃ṭī}, once more
  \item \emph{Say it:} \emph{Darvāze ke pīche ek khū̃ṭī hai.} — There is a peg behind the door.
  \item \emph{Recall:} say the Hindi for a light bulb, then the Hindi for a switch, then say \emph{khū̃ṭī} again
\end{itemize}

\begin{tcolorbox}[breakable,colback=teal!4,colframe=teal!35!black,title={Before you move on}]
What does \hi{खूँटी} mean? (\textquotedblleft{}A peg; a hook (on a wall)\textquotedblright{}.) Say it once more.
\end{tcolorbox}
\section[\texorpdfstring{\hi{टाँगना} (ṭā̃gnā)}{टाँगना (ṭā̃gnā)}]{\hi{टाँगना} (ṭā̃gnā) — to hang up}

\label{lesson:HI-C518-tagna}



\begin{quote}
Before the new one: say the Hindi for a switch, then the Hindi for a peg.
\end{quote}

\subsection*{You'll want to know: \hi{टाँगना}}
\textbf{\hi{टाँगना}} — \emph{ṭā̃gnā} — \textquotedblleft{}to hang up\textquotedblright{}.

\begin{quote}
\hi{मीरा} \hi{ने} \hi{अपना} \hi{छाता} \hi{खूँटी} \hi{पर} \hi{टाँगा।}
\end{quote}

\begin{quote}
\emph{Mīrā ne apnā chātā khū̃ṭī par ṭā̃gā.}
\end{quote}

\begin{quote}
Mira hung her umbrella on the peg.
\end{quote}

Past, with \hi{ने} (he or she …): \hi{उसने} \hi{कोट} \hi{टाँगा}, \emph{usne koṭ ṭā̃gā}.

\begin{cousinweb}
\textbf{\hi{टाँगना}} (\emph{ṭā̃gnā}) is a tadbhava, inherited from \emph{*ṭaṅg}, \textquotedblleft{}hang\textquotedblright{}, a form CDIAL reconstructs from its descendants (CDIAL 5436).
\end{cousinweb}

\begin{grammarlens}[title={What you've built}]
One more word for describing what you lived and felt.
\end{grammarlens}

\subsection*{Your turn}
\begin{itemize}
\raggedright
  \item \emph{Say it:} \emph{ṭā̃gnā}
  \item \emph{Say it:} \emph{ṭā̃gnā}, once more
  \item \emph{Say it:} \emph{Mīrā ne apnā chātā khū̃ṭī par ṭā̃gā.} — Mira hung her umbrella on the peg.
  \item \emph{Recall:} say the Hindi for a switch, then the Hindi for a peg, then say \emph{ṭā̃gnā} again
\end{itemize}

\begin{tcolorbox}[breakable,colback=teal!4,colframe=teal!35!black,title={Before you move on}]
What does \hi{टाँगना} mean? (\textquotedblleft{}To hang up\textquotedblright{}.) Say it once more.
\end{tcolorbox}
